\section{Bias versus purity at power widths}
\label{sec:bias-versus-purity-at-power-widths}

At prescribed power widths, the next lemma combines an available density
surplus with the absence of nearly monochromatic patches at doubled width
budgets to construct a cube. Its contrapositive will be the first reduction
toward discrepancy.

\begin{lemma}[Bias versus purity]\label{lem:bias-purity}
Fix constants \(0<\beta\le\gamma<1\). For \(h=h(\beta,\gamma)>0\) chosen below, the following two hypotheses cannot both hold along a bad sequence:
\begin{enumerate}[label=\textup{(\roman*)}]
\item\label{it:bias-purity-available} Pairs of laws at widths \(\le n^\beta,n^\gamma\) on \(X,Y\) respectively with absolute bias \(\ge n^{-h}\) from half are available;
\item\label{it:bias-purity-excluded} On the current sets there is \textbf{no} pair at respective widths \(\le2n^\beta,2n^\gamma\) having defect \(\le\exp(-n^h)\) in either color. Defect for a color means density of its complement.
\end{enumerate}
Equivalently, their simultaneous validity yields a monochromatic cube.
\end{lemma}

\begin{proof}
We will first prepare pairs of laws \((\mu_i,\nu_i)\) with a density
surplus in one color \(G\), and distribute these pairs over the cube.
An even role will share a hidden tuple of \(X\)-labels with nearby even
roles. The odd roles receive laws on \(Y\) restricted to common neighbors
of the tuples they encounter. After assigning distinct odd labels, the
posterior law of an even role's tuple will show that those labels have a
large common neighborhood in \(X\). Masks on the patch laws will balance
the resulting label loads, allowing Hall's theorem to assign the even
roles distinctly as well.

\noindent\textbf{Step 1: a density plateau.}\quad
Set
\[
 \omega=\min(\beta,1-\gamma)/1000,\quad h=\omega/10^6,\quad
 k=\lceil n^{\omega/3}\rceil,\quad L=n^h,\quad a_*=n^{-h}.
\]
Here \(k\) will be the tuple length, \(e^{-L}\) the minimum acceptable
fraction retained by a single common-neighbor restriction, and \(a_*\)
the density surplus we must preserve.

Take residual label sets containing a witness in
hypothesis~\ref{it:bias-purity-available}. For each integer
\(0\le j\le\lfloor n^{\omega/2}/2\rfloor\), let \(P_j\) be the largest
density in either color among laws on those sets with width budgets
\[
 (s_X(j),s_Y(j))=(n^\beta,n^\gamma)
             +j(n^{\beta-\omega/2},n^{\gamma-\omega/2}).
\]
These maxima exist by compactness and form a nondecreasing sequence in
\([0,1]\), starting at least at \(1/2+a_*\). If every consecutive
increase were at least \(n^{-\omega/3}\), their sum would exceed one
for large \(n\). Choose a step with \(P_{j+1}-P_j<n^{-\omega/3}\),
and a maximizing pair at its smaller budgets. Write these budgets as
\(s_X,s_Y\), its color as \(G\), its first law as \(\mu_i\), and its
density as \(p=1/2+a_i\), where \(a_i\ge a_*\).

We can trim the second law to a law \(\nu_i\) such that
\[
 d_G(\mu_i,y)\ge p-n^{-\omega/5}
       \quad(y\in\operatorname{supp}\nu_i),
 \qquad \operatorname{width}(\nu_i)\le s_Y+1.
\]
To see this, first consider the columns with degree greater than
\(p+n^{-\omega/3}\). If their mass exceeded
\(\eta=\exp(-n^{\gamma-\omega/2}/2)\), conditioning on them would
increase width by less than \(n^{\gamma-\omega/2}/2\) and produce a
density larger than \(P_{j+1}\). Their mass is therefore at most \(\eta\).
If \(t\) is the mass of columns below \(p-n^{-\omega/5}\), the original
mean \(p\) gives
\[
 t\bigl(n^{-\omega/5}+n^{-\omega/3}\bigr)
       \le n^{-\omega/3}+\eta.
\]
Thus \(t=o(1)\), so retaining the other columns costs less than one unit
of width. We also retain the following upper density bound: any laws
supported in this patch with respective widths at most
\[
 s_X+\tfrac12 n^{\beta-\omega/2},\qquad
 s_Y+\tfrac12 n^{\gamma-\omega/2}
\]
have \(G\)-density at most \(p+n^{-\omega/3}\). Both laws fit the next
budget pair, which proves this bound independently of the trimming.
The index \(i\) records the laws, budgets, and density just obtained.

This construction works after every removal allowed by
hypothesis~\ref{it:bias-purity-available}, so the resulting patches are
available. Lemmas~\ref{lem:stabilization} and~\ref{lem:balanced-profiles}
allow us to pass to a subsequence, fix one available color \(G\), and use
a finite balanced mixture of its patches.

\noindent\textbf{Step 2: a coordinate key distributing the patches.}\quad
We need a function \(g\) on cube roles whose values are spread over many
keys, while only a small number of keys occur among the neighbors of any
one role. It will use \(\lfloor n^{\gamma-\omega}\rfloor\) disjoint
coordinate gadgets. In each gadget take \(s'\) chunks of \((s')^6\) bits,
where \(S=s'+1\) is a power of two and \(s'\asymp n^{2\omega}\).
The number of bits used by all the gadgets is
\[
 m=\Theta(n^{\gamma+13\omega})=o(n).
\]

For one gadget, clip each chunk's count of ones to a central interval
consisting of \(S^2\) integer unit steps, and translate that interval to
\([0,S^2]\). Denote the clipped values by \(t_1,\ldots,t_{s'}\), with
chunk labels retained. Each endpoint has probability at least \(1/3\):
the counts are symmetric, and their maximum binomial atom is
\(O((s')^{-3})\), so the unclipped central interval has probability
\(O(S^2(s')^{-3})=o(1)\).

Sort these values by rank and add sentinel ranks \(0,S\) with values
\(0,S^2\). A binary search starts from the rank interval \([0,S]\).
At every step its rank endpoints \(a,b\) have values enclosing
\([aS,bS]\). At the middle rank, compare its count with the fixed value
\((a+b)S/2\). If the count is at least this value, continue in the left
half; otherwise continue in the right half. The invariant is preserved.
At a leaf, two consecutive ranks enclose one elementary interval of
length \(S\). The gadget outputs the index of that interval and the
labeled subset of chunks above its midpoint.

There are two useful properties of this output. First, all one-bit
changes to a fixed input produce only \(O(\log n)\) possible outputs.
If the search path is unchanged, membership above the leaf midpoint is
unchanged. Otherwise consider the first comparison that changes on the
original path. One clipped count must move between a multiple of \(S\)
and the integer immediately below it. Once that comparison and the
direction of the move are specified, the new multiset of counts is fixed.
The identity of a tied chunk making the move does not affect the output:
its count stays on the same side of every elementary midpoint. There are
only \(O(\log n)\) comparisons on the original path and two directions.
Second, for any fixed output, each chunk is required to lie on a specified
side of a fixed midpoint. Either requirement has probability at most
\(2/3\), because the opposite clipped endpoint has probability at least
\(1/3\). Independence of the chunks bounds the probability of that output
by \((2/3)^{s'}\).

Let \(g\) be the joint output of all gadgets, and for odd \(u\) let
\(Z_u=\{g(v):v\sim u\}\). Flipping one bit changes at most one gadget,
so
\[
 |Z_u|\le n^{\gamma-\omega+o(1)},\qquad
 \max_g\Pr(g(v)=g)\le\exp(-\Omega(n^{\gamma+\omega})).
\]
In the second estimate, \(v\) is a uniform role on either fixed cube
parity. The unused coordinates allow the special bits to remain uniform
after parity is fixed.

Assign a patch index \(i_g\) to each key so that the averages of
\(\mu_{i_{g(v)}}\) over \(v\in A\), and of \(\nu_{i_{g(u)}}\) over
\(u\in B\), are pointwise at most \(K/N\) for a fixed \(K\).
Such an assignment follows by drawing the tags independently from the
balanced mixture. For example, if \(\alpha_g\) is the fraction of even
roles with key \(g\), each summand in the normalized mean at a fixed
label is bounded by
\[
 \alpha_g N\mu_{i_g}(x)
 \le \exp\bigl(-\Omega(n^{\gamma+\omega})+2n^\gamma\bigr).
\]
The same bound holds on the odd side. The summand bound is at most
\(\exp(-c n^{\gamma+\omega})\) for some fixed \(c>0\), and the total
mean is bounded. Bounded-summand concentration therefore bounds a fixed
positive deviation from that mean by
\(\exp(-\Omega(\exp(c n^{\gamma+\omega})))\).
Since there are at most \(2N=\exp(O(n))\) label
coordinates, one assignment meets all the bounds. Fix it henceforth.

\subsection{Hidden tuples and height choices}
\label{sec:hidden-tuples-and-height-choices}

\noindent\textbf{Step 3: tuples and the masses retained by their filters.}\quad
The height device will let nearby even roles choose from a common set of
tuple indices. Use Lemma~\ref{lem:height-selection} on the even sites of
\(Q_n\), with
\[
 D=2,\quad b=\omega/30,\quad b_0=b/4,\quad \rho=b_0/10,
 \quad r=\lfloor n^{1-\rho}\rfloor,\quad \lambda=n^{10},
 \quad\zeta=b_0/20,\quad\sigma=\zeta/10.
\]
Choose \(\theta,a\) there to satisfy its inequalities, and use its height
bound \(H\). The logarithm of the volume of an \(O(r+H)\)-neighborhood
is \(o(a_*n)\). Set \(T=\lceil3n^b\rceil\); it bounds the number of
indices that an odd role's neighbors will collectively select.

A mask for a law is a subset of its support of mass at least \(1/2\);
the superscript \(M\) denotes conditioning on that subset. For now give
each mask an arbitrary distribution, independently of all other masks.
For every possible center ID \(c\) (a location and a level) and key \(g\),
choose a mask for \(\mu_{i_g}\) and draw a tuple \(W_{c,g}\) of \(k\)
independent labels from \(\mu_{i_g}^M\). Tuples and masks for different
pairs \((c,g)\) are independent. Each odd role \(u\) also has an
independent mask for \(\nu_{i_{g(u)}}\). We generate these data even for
IDs whose centers will be absent. They are independent of prospective
center positions, activations, and the ties used to choose active centers.

The same tuple \(W_{c,g}\) will be used in every odd filter that mentions
\((c,g)\). If an even role \(v\) selects \(c\), and its odd neighbors
eventually receive labels \(\mathbf y=(y_u:u\sim v)\), its available
label set is
\[
 S_{v,c}(\mathbf y)=
 \{x\in\operatorname{supp}\mu_{i_{g(v)}}^M:
                         xy_u\in G\text{ for every }u\sim v\},
\]
where the mask is the one belonging to \((c,g(v))\). We will bound the
posterior density of \(W_{c,g(v)}\), supported on this set's \(k\)th
power, to prove that the set is large. The next estimates quantify how
much an odd filter changes when this entire shared tuple is removed.

For an odd role \(u\) and an ID set \(\mathcal D\) with
\(1\le|\mathcal D|\le T\), define the retained mass
\[
 R_u(\mathcal D)=\sum_y\nu_{i_{g(u)}}^M(y)
       \prod_{c\in\mathcal D}\prod_{g'\in Z_u}
       \prod_{x\text{ an entry of }W_{c,g'}}1_{\{xy\in G\}}.
\]
Let \(R_u^{-(c,g')}(\mathcal D)\) be the same expression with all
requirements from the single indexed tuple \(W_{c,g'}\) omitted.
Other tuples with ID \(c\) remain in this expression. We call the set
\(\mathcal D\) valid at \(u\) if
\[
 R_u(\mathcal D)\ge\exp(-Lk|\mathcal D||Z_u|)
\]
and, for every \(c\in\mathcal D\) and \(g'\in Z_u\),
\[
 \frac{R_u(\mathcal D)}{R_u^{-(c,g')}(\mathcal D)}\ge
 \begin{cases}
 \exp(k(-\log2+c_1a_*)),&g'=g(u),\\
 \exp(-Lk),&g'\ne g(u),
 \end{cases}
\]
where \(c_1>0\) is a fixed small absolute constant. The first condition
ensures that every denominator in these ratios is positive. These
filters use the common-neighborhood sampling operation of dependent
random choice~\cite[Lemma~2.1]{fox-sudakov2009}. We also need their mass
ratios because deleting one tuple will later remove dependence on a
hypothetical even candidate.

For fixed masks, positions, and \(\mathcal D\), the probability that
\(\mathcal D\) is invalid is at most \(\exp(-n^{\omega/5})\). We prove
this by exposing the tuple entries successively, with any specified
tuple placed last. Call an exposed prefix successful if every hit
fraction in it is at least \(e^{-L}\); no condition is imposed on later
entries. After a successful prefix, let \(\eta\) be the law on \(Y\)
remaining after those hit requirements. Their total width cost is at most
\[
 O(kT|Z_u|L)=o(n^{\gamma-\omega/2}).
\]
For the next entry \(x\), drawn from its masked first law, put
\(q(x)=d_G(x,\eta)\). This is the fraction of current mass retained by
that entry. Its probability of being below \(e^{-L}\) is at most
\(\exp(-n^{\beta-\omega/2}/4)\). Otherwise conditioning the incoming
first law on these values of \(x\) would cost at most
\(n^{\beta-\omega/2}/4\) additional width. Together with \(\eta\), it
would give an opposite-color patch of defect below \(e^{-L}\), within
the doubled budgets, contrary to hypothesis~\ref{it:bias-purity-excluded}.

When this entry has key \(g(u)\), both laws lie in the same prepared
patch \(i=i_{g(u)}\). The intrinsic upper density bound from Step 1 also
gives
\[
 \Pr\{q(x)>p+n^{-\omega/3}\}
       \le\exp(-n^{\beta-\omega/2}/4).
\]
Indeed conditioning on a larger exceptional set would fit the first
width margin, and \(\eta\) fits the second margin. For the lower mean
bound, note that \(2\mu_i-\mu_i^M\) is a probability law of density at
most two relative to \(\mu_i\). Since \(\eta\) is supported on the
trimmed second support, the column lower bound and then the intrinsic
upper bound give
\[
 \begin{aligned}
 \mathbb E_{x\sim\mu_i^M}q(x)
 &=2d_G(\mu_i,\eta)-d_G(2\mu_i-\mu_i^M,\eta)\\
 &\ge2(p-n^{-\omega/5})-(p+n^{-\omega/3})
 \ge p-3n^{-\omega/5}.
 \end{aligned}
\]
These estimates hold conditional on each fixed successful prefix.
Comparing the mean with the upper bound for \(q\), and allowing the
exponentially small exceptional set, shows that
\[
 \Pr\{q<p-a_i/4\mid\text{previous successful entries}\}
       =O(n^{-\omega/5}/a_i)=o(a_*/L).
\]

Call such an entry a dip. Its indicator is set to zero once an earlier
entry has had \(q<e^{-L}\). This stopped sequence is adapted to the
successive exposures, and its conditional expectation bound remains
valid throughout; we do not condition on eventual validity. The bounded-increment
inequality~\cite{azuma1967} shows that an own-key last tuple has at most
\(a_*k/(20L)\) dips, except with probability
\(\exp(-\Omega(ka_*^2/L^2))\). On a path with no step below \(e^{-L}\),
each dip costs at most \(L\) in the negative logarithm of the retained
fraction; at every other own-key entry,
\(\log q\ge-\log2+.7a_i\). Multiplying the \(k\) fractions gives the
own-key ratio above, for sufficiently small \(c_1\). For a cross-key
tuple the lower bound \(q\ge e^{-L}\) already gives its ratio. The same
bound on all entries gives the lower bound for \(R_u(\mathcal D)\).
There are only polynomially many entries and choices of the last tuple,
and \(ka_*^2/L^2=n^{\omega/3-4h+o(1)}\). A union bound therefore gives
the claimed \(\exp(-n^{\omega/5})\) failure estimate.

\noindent\textbf{Step 4: excluding failed sets and selecting heights.}\quad
We now define eligibility before activating centers. For each odd role
\(u\) and consecutive pair of levels, consider all ID sets of size at
most \(T\) from the prospective centers in the balls of its incident
even sites at these levels. Among the sets invalid at \(u\), take a
maximal disjoint family, using a fixed local rule, and mark every ID in
its union forbidden for those even sites, at that ID's own level. Every
invalid set intersects this union; otherwise it could be added to the
family. Consequently every set consisting entirely of the remaining
eligible IDs is valid.

Binomial tails show that, with probability tending to one, every
site-level ball contains between \(\lambda/2\) and \(2\lambda\)
prospective centers. Fix such positions and all masks. For a given odd
role and level pair, a prospective list of at most \(T\) IDs can be
chosen in \(\exp(O(T\log n))\) ways. Invalidity events for disjoint ID
sets depend on disjoint tuples, hence are independent under this
conditioning. The probability of a disjoint family of \(n\) invalid
sets is therefore at most
\[
 \exp\bigl(n[O(T\log n)-n^{\omega/5}]\bigr).
\]
Since \(T=O(n^{\omega/30})\), this bound still tends to zero after the
union over all odd roles and level pairs.

Thus each such family has fewer than \(n\) members, each with at most
\(T\) IDs. An even site has \(n\) neighboring odd roles, so at most
\(O(n^2T)\) IDs are forbidden in any one of its level balls. This is
\(o(\lambda)\), leaving at least \(\lambda/3\) eligible centers.
Activate the centers and apply the long local height rule of
Lemma~\ref{lem:height-selection}, with uniform eligible active choices.
With probability \(1-o(1)\), the stated counts hold and the rule gives
correct heights everywhere. Around an odd role, these heights occupy at
most two consecutive levels because its even neighbors are at distance
two from one another. At each occupied level the crowd bound in a
representative site's \((r+2)\)-ball allows at most \(n^b\) selected
IDs. The actual set therefore has size at most \(2n^b<T\) and, by the
marking argument, is valid. Call this global count, eligibility, and
height outcome \emph{geometric success}.

At an odd role \(u\) with valid incident choices, let \(p_u\) be
\(\nu_{i_{g(u)}}^M\) restricted to the corresponding tuple hits and
normalized. If the choices fail, or their ID set is too large or invalid,
define \(p_u=0\). Thus it is always a subprobability. Auxiliary product
experiments may use fixed fallback samples for these zero rows; the
fallbacks do not replace the zero output in a load estimate. Every rule
uses only data within distance \(O(r+H)\): height choices consult their
long tubes, and marks consult the nearby balls, tuples, and odd masks.
No rule consults geometric success itself. On that success all odd rows
are probabilities, and the retained-mass bound always gives
\[
 Np_u(y)\le\exp(O(n^\gamma)).
\]

\noindent\textbf{Step 5: the posterior tuple and even common neighbors.}\quad
Fix an even role \(v\) and a possible reference ID \(c\) in its ball
at some level. Let \(E\) be the following local event: \(c\) is present
and is legitimately selected at \(v\) by the long path rule and its tie,
at a good height below \(H\); eligibility sizes are correct throughout
\(v\)'s height tube; all neighboring odd kernels are valid; and the
prospective counts in every neighboring-filter ball at all levels are
at most \(2\lambda\). These conditions hold for the selected ID on
geometric success.

\Needspace{6\baselineskip}
Hold all preparatory data except \(z=W_{c,g(v)}\) fixed, including its
mask. Its prior \(\pi\) remains the original product of \(k\) masked
first laws. For the vector \(\mathbf y=(y_u:u\sim v)\) of odd-neighbor
labels, define
\[
 F_z(\mathbf y)=1_E\prod_{u\sim v}p_u(y_u),\qquad
 M_c(\mathbf y)=\int F_z(\mathbf y)\,d\pi(z).
\]
Selection, validity, and all other local rules in this formula are
recomputed for each replacement of \(z\). Thus \(F_z\) is the
likelihood including the event of selecting this tuple, not the
likelihood conditional on its actual selection.

To remove dependence on \(z\) from the reference law, form a mixture
at each neighbor \(u\). Its terms range uniformly over ID sets of size
at most \(T\) containing \(c\), drawn from the present IDs in the
incident balls at all levels. If the prospective count condition fails
or this family is empty (in particular, if \(c\) is absent), \(E\) is
impossible and we use a fixed reference probability instead of the
whole mixture. Otherwise the count condition makes the
logarithm of the number of terms \(O(T\log n)\). In each term omit
\emph{every} requirement from the indexed tuple \(W_{c,g(v)}\), then
normalize the remaining filtered law; use a fixed default if undefined
or empty. These terms do not depend on \(z\). In particular, they do
not require the proposed set to pass its tests or to be selected.
Let \(Q\) be the product of these mixtures over the neighbors of \(v\).

On \(E\), each actual odd kernel compares to its corresponding mixture
term by the mass ratio from Step 3. All but at most \(m\) neighbors have
\(g(u)=g(v)\), since only flips in special coordinates can change the
key. The comparison cost for the product is consequently at most
\[
 \exp\bigl((n-m)k(\log2-c_1a_*)+mkL+O(nT\log n)\bigr).
\]
Here \(T\log n=o(ka_*)\) and \(mL=o(a_*n)\). For a fixed
\(0<c_2<c_1\), these two error terms fit within
\((c_1-c_2)a_*kn\), giving
\[
 F_z(\mathbf y)\le\exp((\log2-c_2a_*)kn)\,Q(\mathbf y).
\]
The entire tuple has been deleted in \(Q\), while true selection and
validity remain in \(F_z\). This is precisely the distinction needed
for Lemma~\ref{lem:gated-posterior}.

Require, for the selected reference, that
\[
 M_c>0,\qquad M_c\ge e^{-c_2a_*kn/4}Q.
\]
The posterior \(F_z\,d\pi(z)/M_c\) then has density relative to the
uniform product law on the original \(N\)-label side \(X\), with current
laws extended by zero to discarded labels, at most
\[
 \exp((\log2-c_2a_*/2)kn).
\]
Indeed division by the displayed threshold uses \(c_2a_*kn/4\) of the
gain, and the masked first width is \(o(a_*n)\), leaving the stated
bound after multiplying the \(k\) prior density factors.
Writing \(S=S_{v,c}(\mathbf y)\), the posterior is supported on
\(S^k\), so
\[
 1\le\exp((\log2-c_2a_*/2)kn)(|S|/N)^k,
 \qquad |S|\ge N\exp(-(\log2-c_2a_*/2)n).
\]
Define \(p_v\) to be uniform on \(S\) if the selected reference meets
\(E\) and the predictive threshold, and zero otherwise. This even
subprobability uses only the local preparatory data and neighboring odd
labels. Whenever it is nonzero, its support is a common neighborhood in
the same color \(G\).

\subsection{Balance and load transfer}
\label{sec:balance-and-load-transfer}

\noindent\textbf{Step 6: choosing mask distributions to balance the rows.}\quad
For this step use the raw experiment with conditionally independent odd
draws given the preparatory data, including fallbacks where necessary,
and without conditioning on geometric success. We will choose the mask
distributions so that
\[
 \mathbb E p_u\le2\nu_{i_{g(u)}}\quad(u\in B),\qquad
 \sum_{v\in A}\mathbb E p_v
       \le K'\sum_{v\in A}\mu_{i_{g(v)}}
\]
coordinatewise, for a fixed \(K'\).

The player choosing the odd mask at \(u\) controls the output vector
\(\mathbb E p_u\). Given nonnegative label prices on \(Y\), retain the
labels with price at most twice the \(\nu_{i_{g(u)}}\)-mean price.
They have mass at least one half. Every nonzero \(p_u\) is supported
there and has mass at most one, so its expected price is at most twice
that mean, regardless of the other mask profiles. Separation therefore
gives a mixture of masks meeting this player's vector bound.

The player for \((c,g)\) instead controls the sum of outputs from even
rows with key \(g\) selecting ID \(c\). Let \(n_{c,g}\) be the number
of even roles with key \(g\) whose radius-\(r\) balls contain the
location of \(c\), regardless of its presence or selection. This count
is deterministic. Put
\[
 w_c=(\lambda/V)
 \begin{cases}
 3/\lambda,&\text{level zero},\\
 \exp(-n^{c_3}),&\text{otherwise},
 \end{cases}
\]
where \(V\) is the radius-\(r\) ball volume and \(c_3>0\) is small and
fixed. The probability that one such row selects \(c\) with event \(E\)
is at most \(w_c\). Presence of \(c\) costs \(\lambda/V\); given
presence, Lemma~\ref{lem:height-selection} bounds its level-zero choice
probability by \(3/\lambda\), and a positive height by
\(\exp(-n^{c_3})\). The local size condition in \(E\) is retained in
this application. The lemma's position-averaged bound is uniform over
the legal eligibility sets, so the bound remains valid for the masks
and tuple-dependent eligibility used here.

Given prices on \(X\), this player can likewise choose a mask whose
prices are at most twice their \(\mu_{i_g}\)-mean. Every even output
attributed to this player is supported in that mask. Multiplying its
price bound by the probability \(w_c\) of selecting \(c\) shows that
the player's expected output vector can be bounded by
\(2n_{c,g}w_c\mu_{i_g}\), again by separation. The output rules for
fixed masks do not depend on their profile probabilities: in particular,
the posterior computations condition on the realized masks. Thus
Lemma~\ref{lem:balanced-profiles} supplies simultaneous profiles for
all players.

For any fixed even role \(v\), summing \(w_c\) over IDs in its range
costs at most \(3\) at level zero and at most
\(H\lambda\exp(-n^{c_3})=o(1)\) over the other levels. Summing the
players' vector bounds proves the even inequality. The fixed patch-tag
assignment from Step 2 now shows that the parity averages of
\(N\mathbb E p_u(y)\) and \(N\mathbb E p_v(x)\) are bounded by
constants at every label.

\noindent\textbf{Step 7: assigning odd labels and controlling even loads.}\quad
First bound odd column sums in the preparatory experiment. The weights
\(Np_u(y)\) have cap \(\exp(O(n^\gamma))\), bounded average means by
Step 6, and independent local inputs for roles separated by more than
\(O(r+H)\). A near set has parity fraction at most
\(\exp(-n\log2+o(a_*n))\). Hence its contribution times \(n\) and
the cap tends to zero. Lemma~\ref{lem:scattered-moments}, followed by
Markov's inequality and the union over labels, bounds every normalized
average odd load by a constant with probability tending to one.
Since \(|B|/N=1/(2C_n)\to0\), the actual column sums are then at most
\(1/10\).

Write \(\mathcal H\) for the preparatory history and
\(S_{\rm pre}\) for geometric success together with the odd column
bound. On \(S_{\rm pre}\), all odd rows are probabilities and meet the hypotheses of
Lemma~\ref{lem:near-product-injection} with \(d=N\): their maximum atom
is \(N^{-1}\exp(O(n^\gamma))\le N^{-.95}\) eventually, the column
bound is less than \(.4\), and \(n^2\le N^{.025}\). Let
\(\mathsf J_{\mathcal H}\) be the resulting odd injection law, and
let \(\mathsf P_{\mathcal H}\) be the auxiliary product law of the odd
draws. For every nonnegative function \(\Phi\) of at most \(n^2\) odd
outputs, the joint target bound gives
\[
 \int\Phi(\mathcal H,\mathbf y)\,d\mathsf J_{\mathcal H}
 \le \exp(N^{-.04}n^2)
       \int\Phi(\mathcal H,\mathbf y)\,d\mathsf P_{\mathcal H}
       \qquad(\mathcal H\in S_{\rm pre}).
\]
The comparison factor is \(1+o(1)\), uniformly in the entering history.

We next show that all selected even references meet the predictive
threshold of Step 5 with probability tending to one. For a fixed even
role \(v\) and possible reference \(c\), let
\(\epsilon=e^{-c_2a_*kn/4}\) and let \(D_c\) be the set of neighboring
data with \(M_c=0\) or \(M_c<\epsilon Q\). Apply the comparison with
\(\Phi=1_E1_{D_c}\), then average the preparatory histories with
\(1_{S_{\rm pre}}\) retained. Removing this last indicator increases
the resulting product-draw integral. Its remaining indicator \(1_E\)
makes the neighboring likelihood exactly \(F_z\), because the odd
kernels are valid on \(E\). With the other preparatory data fixed,
integration over the original masked tuple prior therefore gives
\[
 \sum_{\mathbf y\in D_c}\int F_z(\mathbf y)\,d\pi(z)
   =\sum_{\mathbf y\in D_c}M_c(\mathbf y)
   \le\epsilon\sum_{\mathbf y\in D_c}Q(\mathbf y)
   \le\epsilon.
\]
There are at most \(|A|(H+1)V=\exp(O(n))\) role--ID pairs, whereas
\(a_*kn=n^{1+\omega/3-h+o(1)}\). Thus their total predictive-failure
probability is \(o(1)\). The probability here is under the experiment
that first draws the preparatory history and then, on \(S_{\rm pre}\),
draws the odd injection. Failure to reach \(S_{\rm pre}\) also has
probability \(o(1)\); earlier histories have not been reweighted by the
probability of later success.

Let \(S_{\rm good}\) be \(S_{\rm pre}\) intersected with the event
that every selected even reference passes its predictive requirement.
It remains to control the even column sums after this odd injection.
On \(S_{\rm good}\), the normalized even rows have cap
\(\exp((\log2-c_2a_*/2)n)\). With the same near fraction as above,
\[
 n\exp(-n\log2+o(a_*n))
       \exp((\log2-c_2a_*/2)n)=o(1).
\]
For the separated factors in an \(n\)th moment, their even kernels
consult at most \(n^2\) odd rows in total. The event \(S_{\rm good}\)
also tests references at other roles, so first use
\[
 1_{S_{\rm good}}\prod_{j=1}^{\ell}Np_{v_j}(x)
       \le 1_{S_{\rm pre}}\prod_{j=1}^{\ell}Np_{v_j}(x),
       \qquad \ell\le n.
\]
Each factor on the right still contains its own local validity and
predictive requirements. Conditional on the entering history, the product
of these factors is therefore a function of at most \(n^2\) odd outputs,
to which the displayed injection comparison applies. Only after that
comparison remove \(1_{S_{\rm pre}}\), which restricts the preparatory
history. In the resulting raw experiment the separated kernels depend on
disjoint center data, masks, and odd draw seeds. The patch tags are
already fixed. Their expectations thus factor, and their individual
means are those bounded in Step 6. Explicitly, for separated
\(v_1,\ldots,v_\ell\), \(\ell\le n\), and a fixed label \(x\),
\[
 \mathbb E\!\left[1_{S_{\rm good}}
                  \prod_{j=1}^{\ell}Np_{v_j}(x)\right]
       \le(1+o(1))\prod_{j=1}^{\ell}
                    \mathbb E_{\rm raw}[Np_{v_j}(x)].
\]
This is the separated-product hypothesis of
Lemma~\ref{lem:scattered-moments}; the displayed cap estimate handles
nearby factors.

It follows that, except for probability \(o(1)\) on \(S_{\rm good}\), every even
column sum is at most one. Intersect this event with the preceding
success events. All even rows are now probabilities supported on common
\(G\)-neighbors of the assigned odd labels, so Hall's
theorem~\cite{hall1935} assigns them distinct representatives. The even
images lie in \(X\) and the odd images in the disjoint side \(Y\);
every cube edge has color \(G\). This gives the required contradiction
and proves the assertion, including \(\beta=\gamma\).
\end{proof}
